\SourceLeaf{81}%
a est residuum biquadraticum ipsius $p$, puta $\equiv\alpha^4$, quatuor radices congruentiae $x^4\equiv a$ erunt $\alpha,f\alpha,-\alpha,-f\alpha$, quas inter se incongruas esse facile perspicitur. Hinc patet, si colligantur residua minima positiva biquadratorum $1,16,81,256\dotsfour(p-1)^4$, quaterna semper aequalia fore, ita ut $\tfrac14(p-1)$ residua biquadratica diversa habeantur complexum $A$ formantia. Si residua minima biquadratorum usque ad $(\tfrac12p-\tfrac12)^4$ tantum colliguntur, singula bis aderunt.

\Article{7}
Productum duorum residuorum biquadraticorum manifesto est residuum biquadraticum, sive e multiplicatione duorum numerorum classis $A$ semper prodit productum, cuius residuum minimum positivum ad eandem classem pertinet. Perinde producta numeri ex $B$ in numerum ex $D$, vel numeri ex $C$ in numerum ex $C$, habebunt residua sua minima in $A$.

In $B$ autem cadent residua productorum $A.B$ et $C.D$; in $C$ residua productum $A.C$, $B.B$ et $D.D$; denique in $D$ residua productorum $A.D$ et $B.C$.

Demonstrationes tam obviae sunt, ut sufficiat, unam indicavisse. Sint e. g. $c$ et $d$ numeri ex $C$ et $D$, atque $c\equiv eea$, $d\equiv e^3a'$, denotantibus $a,a'$ numeros ex $A$. Tunc $e^4aa'$ erit residuum biquadraticum, i. e. ipsius residuum minimum ad $A$ referetur: quare quum productum $cd$ fiat $\equiv e.e^4aa'$, illius residuum minimum in $B$ contentum erit.

Simul facile iam diiudicari potest, ad quamnam classem referendum sit productum e pluribus factoribus. Scilicet tribuendo classi $A,B,C,D$ resp. characterem $0,1,2,3$, character producti vel aggregato characterum singulorum factorum aequalis erit, vel eius residuo minimo secundum modulum $4$.

\Article{8}
Operae pretium visum est, hasce propositiones elementares absque adminiculo theoriae residuorum potestatum evolvere, qua in auxilium vocata omnia adhuc multo facilius demonstrare licet.

Sit $g$ radix primitiva pro modulo $p$, i. e. numerus talis, ut in serie potestatum $g,gg,g^3\dotsfour$ nulla ante hanc $g^{p-1}$ unitati secundum modulum $p$ congrua evadat. Tunc residua minima positiva numerorum $1,g,gg,g^3\dotsfour g^{p-2}$ praeter ordinem cum his $1,2,3\dotsfive p-1$ convenient, et in quatuor classes sequenti modo distribuentur:

\SourceLeaf{82}%
\[
\begin{array}{c|l}
\text{ad}&\text{residua minima numerorum}\\\hline
A&1,\ g^4,g^8,g^{12}\dotsfive g^{p-5}\\
B&g,\ g^5,g^9,g^{13}\dotsfive g^{p-4}\\
C&gg,g^6,g^{10},g^{14}\dotsfive g^{p-3}\\
D&g^3,g^7,g^{11},g^{15}\dotsfive g^{p-2}
\end{array}
\]

Hinc omnes propositiones praecedentes sponte demanant.

Ceterum sicuti hic numeri $1,2,3\dotsfour p-1$ in quatuor classes distributi sunt, quarum complexus per $A,B,C,D$ designamus, ita \emph{quemvis} integrum per $p$ non divisibilem, ad normam ipsius residui minimi secundum modulum $p$, alicui harum classium adnumerare licebit.

\Article{9}
Denotabimus per $f$ residuum minimum potestatis $g^{\frac14(p-1)}$ secundum modulum $p$, unde quum fiat $ff\equiv g^{\frac12(p-1)}\equiv-1$ (\emph{Disquis. Arithm.} art. 62), patet, characterem $f$ hic idem significare quod in art. 6. Potestas $g^{\frac14\lambda(p-1)}$ itaque, denotante $\lambda$ integrum positivum, congrua erit secundum modulum $p$ numero $1,f,-1,-f$, prout $\lambda$ formae $4m,4m+1,4m+2,4m+3$ resp., sive prout residuum minimum ipsius $g^\lambda$ in $A,B,C,D$ resp. reperitur. Hinc nanciscimur criterium persimplex ad diiudicandum, ad quam classem numerus datus $h$ per $p$ non divisibilis referendus sit; pertinebit scilicet $h$ ad $A,B,C$ vel $D$, prout potestas $h^{\frac14(p-1)}$ secundum modulum $p$ numero $1,f,-1$ vel $-f$ congrua evadit.

Tamquam corollarium hinc sequitur, $-1$ semper ad classem $A$ referri, quoties $p$ sit formae $8n+1$, ad classem $C$ vero, quoties $p$ sit formae $8n+5$. Demonstratio huius theorematis a theoria residuorum potestatum independens ex iis, quae in \emph{Disquisitionibus Arithmeticis} art. 115, III docuimus, facile adornari potest.

\Article{10}
Quum \emph{omnes} radices primitivae pro modulo $p$ prodeant e residuis potestatum $g^\lambda$, accipiendo pro $\lambda$ omnes numeros ad $p-1$ primos, facile perspicitur, illas inter complexus $B$ et $D$ aequaliter dispertitas fore, basi $g$ semper in $B$ contenta. Quodsi loco numeri $g$ radix alia primitiva e complexu $B$ pro basi accipitur, classificatio eadem manebit; si vero radix primitiva e complexu $D$ tamquam basis adoptatur, classes $B$ et $D$ inter se permutabuntur.

\SourceLeaf{83}%
Si classificatio criterio in art. praec. prolato superstruitur, discrimen inter classes $B$ et $D$ inde pendebit, utram radicem congruentiae $xx\equiv-1\; (\mathrm{mod.}\,p)$ pro numero characteristico $f$ adoptemus.

\Article{11}
Quo facilius disquisitiones subtiliores, quas iam aggressuri sumus, per exempla illustrari possint, constructionem classium pro omnibus modulis infra $100$ hic apponimus. Radicem primitivam pro singulis minimam adoptavimus.

\[
\begin{gathered}
p=5\\
g=2,\ f=2\\[3pt]
\begin{array}{c|l}
A&1\\
B&2\\
C&4\\
D&3
\end{array}
\end{gathered}
\]

\[
\begin{gathered}
p=13\\
g=2,\ f=8\\[3pt]
\begin{array}{c|l}
A&1,3,9\\
B&2,5,6\\
C&4,10,12\\
D&7,8,11
\end{array}
\end{gathered}
\]

\[
\begin{gathered}
p=17\\
g=3,\ f=13\\[3pt]
\begin{array}{c|l}
A&1,4,13,16\\
B&3,5,12,14\\
C&2,8,9,15\\
D&6,7,10,11
\end{array}
\end{gathered}
\]

\[
\begin{gathered}
p=29\\
g=2,\ f=12\\[3pt]
\begin{array}{c|l}
A&1,7,16,20,23,24,25\\
B&2,3,11,14,17,19,21\\
C&4,5,6,9,13,22,28\\
D&8,10,12,15,18,26,27
\end{array}
\end{gathered}
\]

\SourceLeaf{84}%
\[
\begin{gathered}
p=37\\
g=2,\ f=31\\[3pt]
\begin{array}{c|l}
A&1,7,9,10,12,16,26,33,34\\
B&2,14,15,18,20,24,29,31,32\\
C&3,4,11,21,25,27,28,30,36\\
D&5,6,8,13,17,19,22,23,35
\end{array}
\end{gathered}
\]

\[
\begin{gathered}
p=41\\
g=6,\ f=32\\[3pt]
\begin{array}{c|l}
A&1,4,10,16,18,23,25,31,37,40\\
B&6,14,15,17,19,22,24,26,27,35\\
C&2,5,8,9,20,21,32,33,36,39\\
D&3,7,11,12,13,28,29,30,34,38
\end{array}
\end{gathered}
\]

\[
\begin{gathered}
p=53\\
g=2,\ f=30\\[3pt]
\begin{array}{c|l}
A&1,10,13,15,16,24,28,36,42,44,46,47,49\\
B&2,3,19,20,26,30,31,32,35,39,41,45,48\\
C&4,6,7,9,11,17,25,29,37,38,40,43,52\\
D&5,8,12,14,18,21,22,23,27,33,34,50,51
\end{array}
\end{gathered}
\]

\[
\begin{gathered}
p=61\\
g=2,\ f=11\\[3pt]
\begin{array}{c|l}
A&1,9,12,13,15,16,20,22,25,34,42,47,56,57,58\\
B&2,7,18,23,24,26,30,32,33,40,44,50,51,53,55\\
C&3,4,5,14,19,27,36,39,41,45,46,48,49,52,60\\
D&6,8,10,11,17,21,28,29,31,35,37,38,43,54,59
\end{array}
\end{gathered}
\]

\[
\begin{gathered}
p=73\\
g=5,\ f=27\\[3pt]
\begin{array}{c|l}
A&1,2,4,8,9,16,18,32,36,37,41,55,57,64,65,69,71,72\\
B&5,7,10,14,17,20,28,33,34,39,40,45,53,56,59,63,66,68\\
C&3,6,12,19,23,24,25,27,35,38,46,48,49,50,54,61,67,70\\
D&11,13,15,21,22,26,29,30,31,42,43,44,47,51,52,58,60,62
\end{array}
\end{gathered}
\]

\SourceLeaf{85}%
\[
\begin{gathered}
p=89\\
g=3,\ f=34\\[3pt]
\begin{array}{c|l}
A&1,2,4,8,11,16,22,25,32,39,44\mathpunct{.}\,45,\\
&50,57,64,67,73,78,81,85,87,88\\
B&3,6,7,12,14,23,24,28,33,41,43,46,\\
&48,56,61,65,66,75,77,82,83,86\\
C&5,9,10,17,18,20,21,34,36,40,42,47,\\
&49,53,55,68,69,71,72,79,80,84\\
D&13,15\mathpunct{.}\,19,26,27,29\mathpunct{.}\,30,31,35,37,38\mathpunct{.}\,51,\\
&52,54,58,59,60,62,63,70,74,76
\end{array}
\end{gathered}
\]

\[
\begin{gathered}
p=97\\
g=5,\ f=22\\[3pt]
\begin{array}{c|l}
A&1,4,6,9,16,22,24,33,35,36,43,47,\\
&50,54,61,62,64,73,75,81,88,91,93,96\\
B&5,13,14,17,19,20,21,23,29,30,41,45,\\
&52,56,67,68,74,76,77,78,80,83,84,92\\
C&2,3,8,11,12,18,25,27,31,32,44,48,\\
&49,53,65\mathpunct{.}\,66,70,72,79,85,86,89,94,95\\
D&7,10,15,26,28,34,37,38,39\mathpunct{.}\,40,42,46,\\
&51,55,57,58,59,60,63,69,71,82,87,90
\end{array}
\end{gathered}
\]

\Article{12}
Quum numerus $2$ sit residuum quadraticum omnium numerorum primorum formae $8n+1$, non-residuum vero omnium formae $8n+5$, pro modulis primis formae prioris $2$ in classe $A$ vel $C$, pro modulis formae posterioris in classe $B$ vel $D$ invenietur. Quum discrimen inter classes $B$ et $D$ non sit essentiale, quippe quod tantummodo ab electione numeri $f$ pendet, modulos formae $8n+5$ aliquantisper seponemus. Modulos formae $8n+1$ autem \emph{inductioni} subiiciendo, invenimus $2$ pertinere ad $A$ pro $p=73,89,113,233,257,281,337,353$ etc.; contra $2$ pertinere ad $C$ pro $p=17,41,97,137,193,241,313,401,409,433,449,457$ etc.

Ceterum quum pro modulo primo formae $8n+1$ numerus $-1$ sit residuum biquadraticum, patet, $-2$ semper cum $+2$ ad eandem classem referendum esse.

